\documentclass[11pt]{article}
\usepackage{series}

\title{The Proto-Spinor:\\
Conditional Spinorial Closure and the q79 Interface\\
\large Corrected sixth edition}
\author{Peter Nero}
\date{July 2026}

\begin{document}
\maketitle

\begin{abstract}
We isolate the part of the proto-spinor construction that is mathematically
forced after its geometric premises are stated.  A world-in-world field is a
section of $\operatorname{Hom}(TP,TI)$ for oriented rank-three bundles; its
nine local components do not multiply manifold dimensions.  If closure data
must retain the nontrivial loop class in $SO(3)$, the faithful continuous
double-cover carrier is spinorial and is represented through
$\operatorname{Spin}(3)\cong SU(2)$.  This is a conditional lifting theorem,
not a derivation of internal dimension three or of physical fermions from
admissibility alone.  For the selected q79 degree-three spectral cover we
record the exact trace split of ranks $1+2$, the common $1+2+3$ carrier, and
the local binary-dihedral lift of signed sheet monodromy. The shared circle is
now a specified pullback of one universal flat $\mathbb Z_{64}$ differential
line, not merely an isomorphic fiber. Its determinant, CLN, root-plane, and
finite-Hessian actions intertwine exactly. The BHT and Hori fiber transforms
share one orthogonal complex structure $J_{\rm FM}$, and the binary sheet group
has exactly two conjugate $\pm i$ phase-root lifts for which
$(\pm iJ_{\rm FM})^2=I$. Strict global Spin closure remains equivalent to a
branch-relator obstruction calculation; the closed result is SpinC and finite
symbolic. Circle, lens, and nil remain typed phase, finite-transport, and
anchoring roles. Lorentzian spinors, masses, and particle assignments remain
downstream realization problems.
\end{abstract}

\section*{Revision note for this edition}
\addcontentsline{toc}{section}{Revision note for this edition}
\begin{description}
\item[Supersedes.] \emph{The Proto-Spinor: Conditional Spinorial Closure and
the q79 Interface}, version~5.
\item[Reason.] Version~5 stopped at a local $\operatorname{Dic}_3$ lift and an
untyped shared line. It did not include the later universal differential-line
pullback, the common BHT--Hori Clifford operator, or the exact binary SpinC
return.
\item[Resolution.] Version~6 retains the conditional double-cover theorem and
adds the flat q79 shared-line theorem, finite Hessian naturality, the exact
$J_{\rm FM}^2=-I$ polarization, and the two conjugate $\pm i$ phase-root lifts.
\item[Retained result.] Spinorial lifting is necessary when an oriented
rank-three carrier must retain the nontrivial $SO(3)$ loop class.
\item[Remaining boundary.] Strict global q79 Spin, the physical nonflat
FM/HYM lift, the strain-to-q79 continuum intertwiner, and a Lorentzian
particle/action completion remain open gates.
\end{description}

\tableofcontents

\section{Status and dependencies}

This paper is a conditional carrier theorem.  Its conclusions depend on the
following data, none of which follows from a bare projection map:
\begin{enumerate}
\item oriented Euclidean rank-three bundles $TP$ and $TI$ over a declared
base $B$;
\item a comparison field $Q_{\rm WW}\in
\Gamma(\operatorname{Hom}(TP,TI))$;
\item an $SO(3)$ frame or holonomy bundle whose nontrivial loop class is to be
retained by the state description;
\item for the q79 specialization, the selected degree-three spectral cover,
its trace map, and its sheet monodromy; and
\item for a physical fermion interpretation, a separate Lorentzian spin or
$\operatorname{Spin}^c$ completion, action, observable map, and normalization.
\end{enumerate}

The terms ``forced'' and ``necessary'' below always mean necessary relative to
these premises.  They do not mean that MTT has selected the premises uniquely.

\section{World-in-world comparison data}

\begin{definition}[Comparison field]
Let $TP,TI\to B$ be oriented Euclidean vector bundles of rank three.  A
world-in-world comparison field is
\[
 Q_{\rm WW}\in\Gamma\!\left(\operatorname{Hom}(TP,TI)\right).
\]
At $b\in B$, choices of oriented orthonormal frames represent
$Q_{\rm WW}(b)$ by a matrix in $\operatorname{Mat}(3,\mathbb R)$.
\end{definition}

\begin{proposition}[Dimension and component bookkeeping]
The total space of $TI\to B$ has dimension $\dim B+3$.  If $\dim B=3$, its
dimension is six.  The number nine is instead the fiber dimension of
$\operatorname{Hom}(TP,TI)$.
\end{proposition}

\begin{proof}
A rank-$r$ vector bundle over a $d$-dimensional base is locally
$U\times\mathbb R^r$ and therefore has total dimension $d+r$.
Furthermore, $\operatorname{Hom}(\mathbb R^3,\mathbb R^3)$ has dimension nine.
\end{proof}

At an invertible, orientation-preserving background, polar decomposition gives
$Q_{\rm WW}=RU$ with $R\in SO(3)$ and $U$ symmetric positive definite.  Its
linearized carrier has the orthogonal Frobenius decomposition
\[
 \operatorname{Mat}(3,\mathbb R)
 =\mathfrak{so}(3)\oplus\operatorname{Sym}(3,\mathbb R),
 \qquad 9=3+6.
\]
After a flag is selected, the symmetric part further decomposes as
\[
 \operatorname{Sym}(3,\mathbb R)
 =\mathbb RI_3\oplus\mathcal D_0\oplus\mathcal O,
 \qquad 6=1+2+3,
\]
where $\mathcal D_0$ is traceless diagonal and $\mathcal O$ is symmetric
off-diagonal.  The latter split depends on the selected flag; only the
$1+5$ trace/traceless split is invariant under the full orthogonal group.

Consequently
\[
 1+3\times3=(1+3)+(1+2+3)=4+6=10
\]
is an exact component identity once an ordering scalar is added.  It is not a
proof of a ten-dimensional manifold, a $3+1$ spacetime, or Lorentzian
signature.

\section{Conditional spinorial lifting theorem}

Let $F^+(TI)\to B$ be the oriented orthonormal frame bundle.  Its structure
group is $SO(3)$ and
\[
 \pi_1(SO(3))\cong\mathbb Z_2,
 \qquad
 1\longrightarrow\{\pm1\}\longrightarrow
 \operatorname{Spin}(3)\longrightarrow SO(3)\longrightarrow1.
\]

\begin{definition}[Neutrality with loop memory]
A carrier has neutrality with loop memory when its projected $SO(3)$ frame
returns after a closed orientation loop while its state representation still
distinguishes the two homotopy classes of such loops.
\end{definition}

\begin{theorem}[Conditional double-cover necessity]
Suppose a continuous state carrier over the $SO(3)$ frame data is required to
distinguish the generator of $\pi_1(SO(3))$ and to compose consistently under
concatenation of lifted paths.  Then an ordinary vector representation of
$SO(3)$ is insufficient.  The minimal connected covering group on which the
nontrivial class remains visible is the universal double cover
$\operatorname{Spin}(3)\cong SU(2)$.
\end{theorem}

\begin{proof}
Every $SO(3)$ representation identifies a loop with its endpoint group
element, so both homotopy classes of loops ending at the identity act as the
identity.  Path lifting to the universal cover assigns the two classes the
endpoints $+1$ and $-1$.  Since $\pi_1(SO(3))=\mathbb Z_2$, the universal
connected cover is two-sheeted and no smaller connected cover can retain that
class.  The universal cover is $\operatorname{Spin}(3)\cong SU(2)$.
\end{proof}

\begin{definition}[Proto-spinor]
Relative to the preceding hypotheses, a proto-spinor is a section of an
associated complex rank-two bundle
\[
 S_I=P_{\rm Spin}\times_{SU(2)}\mathbb C^2,
\]
where $P_{\rm Spin}$ is a chosen lift of the relevant oriented frame or
holonomy data.  The word ``proto'' records that no Lorentzian Dirac operator,
statistics law, or particle ontology has yet been supplied.
\end{definition}

The theorem does not prove that an internal world must be three-dimensional.
For example, $SO(2)$ already has nontrivial winding, while higher-dimensional
orientation groups also possess spin covers.  Rank three is a selected
realization premise supported here by the comparison carrier and the q79
degree-three data, not a universal dimensional-minimality theorem.

\section{Selected q79 trace-split carrier}

Let $\pi:C\to B$ be the selected degree-three spectral cover and set
\[
 \mathcal A=\pi_*\mathcal O_C.
\]
The unit and trace define
\[
 p_{\rm cen}=\frac13\,\mathbf1\circ\operatorname{Tr},
 \qquad
 \mathcal A_0=\ker\operatorname{Tr}.
\]

\begin{theorem}[Trace-split rank flag]
If $\mathcal A$ is finite locally free of rank three and the ground field has
characteristic different from three, then
\[
 \mathcal A\cong\mathcal O\oplus\mathcal A_0,
 \qquad
 \operatorname{rank}(\mathcal O,\mathcal A_0,\mathcal A)=(1,2,3).
\]
Consequently the shared-line carrier
\[
 \mathcal H_{\rm CLN}
 =L_{\rm shared}\otimes
 (\mathcal O\oplus\mathcal A_0\oplus\mathcal A)
\]
has rank six.
\end{theorem}

\begin{proof}
The normalized trace projector is idempotent, has image the unit line, and has
kernel $\mathcal A_0$.  Rank additivity gives $3=1+2$, and the displayed
direct sum has rank $1+2+3=6$.
\end{proof}

This theorem selects a direct trace/trace-zero/full carrier.  On a connected
cover with transitive monodromy it does not produce a global ordering of three
individual sheets.

\section{Universal shared differential line}

Let $\mathcal L_{64}^{\rm univ}\to B_\nabla\mathbb Z_{64}$ be the universal
flat line for the primitive character
$\chi_1(n)=\exp(2\pi\mathrm i n/64)$. The unique nontrivial homomorphism
\[
 h_{S_3}:S_3\longrightarrow\mathbb Z_{64},
 \qquad h_{S_3}(\sigma)=32\,\epsilon(\sigma),
\]
sends even permutations to zero and odd permutations to $32$. Pulling
$\mathcal L_{64}^{\rm univ}$ back along the q79 sheet classifying map gives
$(L_{\rm sh},\nabla_{\rm sh})$. For both admissible odd roots
$r\in\{1,33\}$,
\[
 \chi_r\circ h_{S_3}=\operatorname{sgn}.
\]
Hence $L_{\rm sh}$ and the SpinC determinant sign line are canonically
parallel-isomorphic. This equality preserves connection and holonomy; it is
stronger than recognizing the same order-two phase pattern.

\begin{theorem}[Shared-line finite naturality]
The line $L_{\rm sh}$ tensors the trace/trace-zero/full carrier and acts by
scalar holonomy. It therefore commutes with all three lane projectors, the
root-plane quarter-turn $J_{DE}$, and the normalized finite Hessian
\[
 H_{\rm fin}=\kappa_{\rm fin}(I-P_{\rm Haar}),
 \qquad P_{\rm Haar}=\frac1{6}\sum_{g\in S_3}\rho(g).
\]
On the selected two-copy sheet/edge representation,
\[
 \operatorname{rank}P_{\rm Haar}=2,
 \quad
 \operatorname{spec}(H_{\rm fin}/\kappa_{\rm fin})
 =\{0^{\times2},1^{\times4}\},
 \quad H_{\rm TT}=\kappa_{\rm fin}I_2.
\]
Thus the determinant, CLN, root-plane, and finite-Hessian occurrences are
coherent pullbacks of one differential line, with no new fit parameter.
\end{theorem}

\begin{proof}
The determinant statement is the displayed character identity. Scalar
holonomy commutes with every sheet operator. Haar averaging is an orthogonal
projector, and direct decomposition of the selected representation gives the
rank and spectrum. Tensoring by $L_{\rm sh}$ preserves all identities.
\end{proof}

This theorem is flat and finite-symbol exact. The full physical HYM carrier
has nonzero characteristic curvature and cannot literally have this flat
connection. Its promotion requires a spectral-symbol functor and a parallel
continuum Hessian comparison.

\section{Signed sheet monodromy and its lift}

For a permutation $\sigma\in S_3$, let $P_\sigma$ be its permutation matrix
and define
\[
 \rho_+(\sigma)=\operatorname{sgn}(\sigma)P_\sigma.
\]
Since
$\det(\operatorname{sgn}(\sigma)P_\sigma)=1$, this is a representation
$S_3\to SO(3)$.

\begin{theorem}[Local binary-dihedral lift]
The inverse image of $\rho_+(S_3)$ in $\operatorname{Spin}(3)$ is the
binary-dihedral group $\operatorname{Dic}_3$ of order twelve.  Lifts $q_1,q_2$
of adjacent transpositions may be chosen so that
\[
 q_1^2=q_2^2=-1,
 \qquad q_1q_2q_1=q_2q_1q_2,
 \qquad (q_1q_2)^3=-1.
\]
\end{theorem}

\begin{proof}
The signed permutation image is the rotational symmetry group of an oriented
triangular frame, isomorphic to the dihedral rotation group of order six.  The
inverse image of a finite rotation group under the two-to-one map
$SU(2)\to SO(3)$ has twice its order and is its binary counterpart.  Standard
half-angle quaternion lifts of the two generating rotations satisfy the
displayed relations, identifying the inverse image with
$\operatorname{Dic}_3$.
\end{proof}

\subsection{Global obstruction}

The preceding theorem is not yet a strict global q79 Spin theorem.  Let
$B^\circ$ be the branch complement and let
$\rho:\pi_1(B^\circ)\to SO(3)$ be the signed sheet monodromy.  A strict lift is
a homomorphism $\widetilde\rho$ making
\[
\begin{array}{ccc}
 &\operatorname{Spin}(3)&\\[-2mm]
 &\downarrow&\\[-1mm]
 \pi_1(B^\circ)&\xrightarrow{\rho}&SO(3)
\end{array}
\]
commute.  It exists exactly when the pulled-back central extension is trivial,
equivalently when the relevant $w_2$ obstruction vanishes.  Computationally,
one must lift a presentation of $\pi_1(B^\circ)$ and verify that every relator
closes with central sign $+1$.  Local braid relations alone do not decide all
global relators.

A shared $U(1)$ line can instead participate in a
$\operatorname{Spin}^c$ cancellation of the mod-two obstruction.  That is a
distinct theorem and must not be reported as a strict Spin lift.

\section{Clifford polarization and exact SpinC return}

On the oriented two-circle exterior algebra with basis $(1,p,q,pq)$, the
Poincare-kernel transform is the integral operator
\[
 J_{\rm FM}=
 \begin{pmatrix}
 0&0&0&-1\\
 0&0&-1&0\\
 0&1&0&0\\
 1&0&0&0
 \end{pmatrix}.
\]
Direct multiplication gives
\[
 J_{\rm FM}^{T}J_{\rm FM}=I,
 \qquad \det J_{\rm FM}=1,
 \qquad J_{\rm FM}^{2}=-I.
\]
Thus $J_{\rm FM}$ is an orthogonal complex structure, equivalently one
$\mathrm{Cl}_{0,1}$ generator, and
\[
 P_\pm=\frac12(I\mp\mathrm iJ_{\rm FM})
\]
are complementary complex rank-two polarizations. The same derived operator
reproduces the selected BHT and two-circle Hori fiber transforms.

The half-turn character of $S_3$ cannot itself cancel $J_{\rm FM}^{2}$:
an involution maps to $0$ or $32$ in $\mathbb Z_{64}$, while the square roots
of $32$ are $16$ and $48$. The correct domain is the binary sheet group
$\operatorname{Dic}_3$ already obtained above.

\begin{theorem}[Binary-sheet SpinC--Fourier return]
There are exactly two homomorphisms
\[
 \widetilde h_\pm:\operatorname{Dic}_3\longrightarrow\mathbb Z_{64}
\]
which send the binary central element to $32$ and the two braid generators to
$(16,16)$ or $(48,48)$. Under either odd character these phases are $+i$ and
$-i$, respectively. In the Fourier frame the binary transposition lift is
$J_{\rm FM}$, and therefore
\[
 U_\pm=\pm\mathrm iJ_{\rm FM},
 \qquad U_\pm^2=I.
\]
The two lifts are complex conjugates with the same determinant sign line. No
continuous or discrete fitted parameter is introduced.
\end{theorem}

\begin{proof}
The binary relations force both generator images to be equal square roots of
$32$; exhaustive solution in $\mathbb Z_{64}$ leaves only $16$ and $48$.
The real spinor representation identifies the binary transposition with
$J_{\rm FM}$. Multiplication then gives
$(\pm\mathrm iJ_{\rm FM})^2=I$.
\end{proof}

This closes the double return at representation and flat differential-line
tier. It does not identify compact phase with physical time, choose one of the
conjugate orientations as the observed universe, or supply the nonflat
physical FM/HYM correspondence.

\section{Circle, lens, and nil as typed carrier roles}

In this revision the three labels have the following nonexhaustive roles:
\begin{description}
\item[Circle.] The common line $L_{\rm shared}$ stores phase or holonomy data.
It is counted once and is not physical time.
\item[Lens.] Finite quotient, signed-sheet, or projective transport data.  This
does not require the global compactification to contain a literal lens-space
factor.
\item[Nil.] Anchoring, upper-triangular transport, or termination data in a
selected operator filtration.  This does not prove a literal Nil$_3$ factor.
\end{description}

Thus circle--lens--nil is a carrier and obstruction taxonomy.  The auxiliary
model $L(3,1)\times\mathrm{Nil}_3$ may realize some roles, but it is not the
selected q79 Fu--Yau compactification and the two manifolds are not identified.

Boothby--Wang geometry gives a precise supporting model for this typing. A
prequantum circle bundle over $\mathbb{CP}^1$ with first Chern number $k$ has
total space $L(k,1)$, while an integral circle bundle over $T^2$ has a
Heisenberg nilmanifold as total space. Thus $L(3,1)$ and Nil$_3$ are parallel
circle-bundle realizations over different bases and curvature classes, not
nested spaces \cite{BoothbyWang,ContactBlowup}. Their curved connections are
not the flat q79 root-stack connection. All three may map to the universal
classifier $BU(1)_\nabla$, but an MTT identification requires selected
classifying maps and coherent comparison data.

\section{Downstream physical gates}

\begin{enumerate}
\item A Lorentzian spinor requires a globally hyperbolic base, a spin or
$\operatorname{Spin}^c$ structure, a tetrad, a compatible connection, and a
hyperbolic Dirac operator.
\item A particle interpretation requires a physical state space, statistics,
observables, interactions, and an asymptotic or local excitation criterion.
\item A mass is read from a canonically normalized quadratic action or a
propagator pole.  A positive closure cost alone is not a mass.
\item A Higgs identification requires a selected scalar representation and
couplings.  A positive Hessian does not imply a unique scalar mode.
\item A discrete spectrum requires a self-adjoint operator with compact
resolvent, confining boundary conditions, or another explicit spectral
compactness theorem.  Nil language alone does not quantize a theory.
\end{enumerate}

Weyl, Dirac, Majorana, and twistor descriptions may be built as conditional
encodings once their respective chirality, pairing, real-structure, and
conformal hypotheses are supplied.  They are not consequences of the
rank-three comparison field alone.

\section{The same-source interface theorem still required}

The local strain carrier and the q79 trace-split carrier both have dimensions
$1+2+3$.  Rank equality supplies only a fiberwise linear isomorphism.  The
needed physical theorem must construct
\[
 \mathfrak I_{\rm WW\to q79}:
 \mathbb RI_3\oplus\mathcal D_0\oplus\mathcal O
 \longrightarrow
 L_{\rm shared}\otimes
 (\mathcal O\oplus\mathcal A_0\oplus\mathcal A)
\]
on the selected base and prove:
\begin{enumerate}
\item compatibility with transition functions and the selected flag;
\item isometry for the declared strain and HYM inner products;
\item covariant compatibility
$\mathfrak I\nabla^{\rm WW}=\nabla^{q79}\mathfrak I$; and
\item intertwining of the Hessian, retarded, and overlap operators used by the
numerical source packets.
\end{enumerate}
Until these clauses are proved, the local and q79 objects are related by a
dimension-matched conditional bridge, not an identity. The universal-line and
finite SpinC/Fourier theorems above close important internal squares on the q79
side; they do not supply this world-in-world-to-continuum map.

\section{Relation to the numerical SM program}

The current calculation repositories close embedded renormalized-Standard-
Model equivalence at the declared one-shared-physical-primitive/profile
standard, including the finite matrix, charged Yukawa profile, electroweak,
threshold, and precision transport packets.  Those are valid at their stated
certificate tier.  They do not prove that the proto-spinor carrier uniquely
selects the measured profile, and they do not close the strict no-knob or
unique-observed-branch upgrade.  This paper therefore uses the calculations as
downstream realization evidence, not as a proof of its upstream premises.

\section{Conclusion}

The corrected proto-spinor result is narrower and stronger.  Under explicit
rank-three orientation and loop-memory premises, the double cover is indeed
spinorial.  The selected q79 carrier has an exact $1+2+3$ trace split and an
exact universal flat differential line, an exact $\operatorname{Dic}_3$ lift,
and an exact SpinC--Fourier return on its finite carrier. What remains is
sharply separated: the strict global Spin obstruction calculation, the
physical nonflat FM/HYM lift, the same-source local-to-q79 continuum
intertwiner, and the Lorentzian action/particle completion.

\begin{thebibliography}{9}
\bibitem{LawsonMichelsohn}
H.~B.~Lawson and M.-L.~Michelsohn,
\emph{Spin Geometry}, Princeton University Press, 1989.

\bibitem{FuYau}
J.-X.~Fu and S.-T.~Yau,
\emph{The theory of superstring with flux on non-Kahler manifolds and the
complex Monge--Ampere equation}, J. Differential Geom. 78 (2008).

\bibitem{MTTq79}
P.~Nero,
\emph{Selected q79 Trace-Split CLN Carrier and World-in-World Bridge},
internal theorem packet, 2026.

\bibitem{BoothbyWang}
W.~M.~Boothby and H.~C.~Wang,
\emph{On contact manifolds}, Ann. Math. 68 (1958), 721--734.

\bibitem{ContactBlowup}
R.~Casals, D.~M.~Pancholi, and F.~Presas,
\emph{Contact blow-up}, Expo. Math. 33 (2015), 97--123.

\bibitem{SharedLine}
P.~Nero,
\emph{q79 Universal Shared Differential Line and Finite-Operator
Intertwiner}, executable theorem packet, 2026.

\bibitem{SpinCReturn}
P.~Nero,
\emph{q79 Binary-Sheet FM Shared Root and SpinC Return}, executable theorem
packet, 2026.
\end{thebibliography}

\end{document}
